\section*{Tohu va-bohu}

\epigraf{I la terra era desordenada i buida, i les tenebres planaven sobre la faç de l’abisme.}{Gènesi 1:2}

\begin{quoting}
\texthebrew{וְהָיְתָה הָאָרֶץ תֹהוּ וָבֹהוּ וְחֹשֶׁךְ עַל־פְּנֵי תְהוֹם}

\vspace{1em}

\textit{Veha’aretz hayetah tohu va-vohu vechoshech al-penei tehom}

\vspace{1em}

\textbf{“I la terra era desordenada i buida, i les tenebres estaven sobre la faç de l’abisme.”}
\end{quoting}

\medskip

\noindent
La locució \textbf{\textit{tohu va-bohu}} (\texthebrew{תֹהוּ וָבֹהוּ}) designa l’estat primigeni del món abans de tota creació: un espai sense forma, caòtic, confús.

\begin{itemize}
    \item \textbf{Tohu (\texthebrew{תֹּהוּ}):} desolació, confusió, caos, matèria informe.
    \item \textbf{Bohu (\texthebrew{בֹּהוּ}):} buit inhabitual, buit fosc, absència d’estructura.
\end{itemize}

\noindent
Aquest binomi no designa el no-res absolut, sinó una mena de \textit{plasma caòtic ontològic}, un estat de potencialitat encara no ordenada per la paraula creadora. És el \textit{no-ordre} sobre el qual actua la primera llum.

\bigskip

\noindent
En termes cosmològics, \textit{tohu va-bohu} pot entendre’s com una metàfora del caos informacional previ a la taxonomia, com l’estat d’un núvol digital ple de fragments, arxius duplicats, carpetes sense nom, i dades sense context.